\chapter*{Preface}
\addcontentsline{toc}{chapter}{Preface}

A programming language can be studied by reconstructing a fragment of it in another calculus. The reconstruction begins with a choice about what the fragment does and which aspects of its execution must survive translation. For a stack language, the final stack may suffice to compare arithmetic programs. Once the language has named definitions, a reconstruction may also need to preserve which definition each call resolved to. For a language that touches its environment, the order and results of those interactions may be part of the semantics. Each prototype in this book therefore begins with an account of what a program can do, proceeds to an implementation, and ends by checking what the implementation preserves.

\SP{} gives these reconstructions a common form. Execution produces a history from which a current state can be recovered. A stack, an environment, a program store, a search path, or a set of installed packages is a view of that history. Two programs can then be compared at the level of their results while the evidence needed to distinguish their executions is retained. The claim that one language reconstructs another acquires a concrete meaning: the translation can be exhibited and run, and the observations under which it agrees with the source language can be stated and tested.

The book has three parts. Part~I fixes a small reference kernel with four primitives, states the method used for every later prototype, and works the method on a four-bit shift register, a machine small enough that every transition can be checked. Part~II builds working prototypes: a Forth, a lexically scoped Lisp, a Racket-style module layer with require and provide, domain-specific languages made from module languages and macros, Commodore-style BASIC with its program store and disk \texttt{LOAD}, several shell dialects treated as loading languages, and package indexes with flat and nested resolution. Part~III shows that the prototypes can share one ledger without disturbing each other, compares their structural profiles, and states open problems.

A thread runs through Part~II and explains its order. Racket, BASIC's \texttt{LOAD}, the shell's \texttt{PATH}, and a package manager are four answers to one question: given a name, which of several ordered places supplies the thing it denotes, and what happens to the program once it is found? The prototypes make that shared structure explicit, because each of them implements it with the same ordered-search routine and records every probe as a sphere. A miss becomes a documented refusal and the hit becomes a committed pop.

The method extends an earlier treatment of a Forth stack \cite{flyxion-forth} and uses the four primitives as specified in the \SP{} specification \cite{flyxion-spec}.

Every listing in this book is cut from source files that accompany it, and every transcript is the output of a script that runs those files. The accompanying test suite contains seventy-six tests. Statements about behavior in Chapters~4 through 10 refer to the prototypes, and several details of real systems are deliberately simplified; each chapter says where. Where a statement concerns a real system, such as the order in which a shell resolves a command name, the prototype is a model of the documented rule and not a copy of any implementation.

Notation is fixed once. The four primitives are $\Pop$, $\Refuse$, $\Bind$, and $\Collapse$; a world is a pair $(H,\Omega)$ of a history and a family of live spheres; and a language fragment is written by the name of its tag prefix, such as \texttt{forth.out}.
